\begin{tikzpicture}[
    x=1cm,
    y=1cm,
    every node/.style={font=\small}
]

    % Uložený vzor
    \hopfieldpattern
        {0}{0}
        {\patternThree}
        {uložený vzor\\\(\mathbf{x}_3\)}
        {green}

    % Poškozený vstup
    \hopfieldpattern
        {2.4}{0}
        {\patternNoisy}
        {vstup\\\(\mathbf{s}^{(0)}\)}
        {orange}

    % První mezistav
    \hopfieldpattern
        {4.8}{0}
        {\patternIntermediateA}
        {stav\\\(\mathbf{s}^{(1)}\)}
        {blue}

    % Druhý mezistav
    \hopfieldpattern
        {7.2}{0}
        {\patternIntermediateB}
        {stav\\\(\mathbf{s}^{(2)}\)}
        {blue}

    % Výsledný stabilní stav
    \hopfieldpattern
        {9.6}{0}
        {\patternThree}
        {stabilní stav\\\(\mathbf{s}^{(3)}=\mathbf{x}_3\)}
        {green}

    % Oddělení uloženého vzoru od procesu vybavování
    \draw[dashed,thick,black!45]
        (1.95,-0.15) -- (1.95,2.55);

    % Aktualizace stavů
    \draw[diredge]
        (3.98,0.9) -- (4.72,0.9);

    \draw[diredge]
        (6.38,0.9) -- (7.12,0.9);

    \draw[diredge]
        (8.78,0.9) -- (9.52,0.9);

    % Popisky
    \node[vertexplain,font=\scriptsize]
        at (0.75,-0.22)
        {uložený atraktor};

    \node[vertexplain,font=\scriptsize]
        at (3.15,-0.22)
        {\(30\) chybných bitů};

    \node[vertexplain,font=\scriptsize]
        at (5.55,-0.22)
        {\(8\) chybných bitů};

    \node[vertexplain,font=\scriptsize]
        at (7.95,-0.22)
        {\(2\) chybné bity};

    \node[vertexplain,font=\scriptsize]
        at (10.35,-0.22)
        {žádný chybný bit};

\end{tikzpicture}
